\section{Metodología, diseño y desarrollo}
\label{sec:metodologia}

\subsection{Dataset}
\label{subsec:dataset}
Utilizamos el subconjunto de PlantVillage~\cite{hughes2015plantvillage} correspondiente a
pimiento, papa y tomate: 20\,638 imágenes RGB de $256\times256$ píxeles distribuidas en 15
clases (12 enfermedades y 3 clases de hoja sana). Las imágenes fueron tomadas en condiciones
de laboratorio, con una única hoja centrada sobre un fondo uniforme. El dataset está
desbalanceado (Fig.~\ref{fig:datos}a): la clase mayoritaria (virus del rizado amarillo del
tomate, TYLCV) tiene 3\,208 imágenes y la minoritaria (papa sana) 152, una relación de 21:1.
Las hojas sanas representan el 15,6\,\% del total. Por este desbalance adoptamos el macro-F1
como métrica principal.

\begin{figure}[t]
\centering
\includegraphics[width=\columnwidth]{metodologia_datos.pdf}
\caption{(a) Cantidad de imágenes por clase en el subconjunto de PlantVillage utilizado.
(b) Una hoja de entrenamiento (tizón temprano del tomate) y siete muestras generadas por el
pipeline de data augmentation usado en el fine-tuning ($160\times160$).}
\label{fig:datos}
\end{figure}

\subsection{Partición de datos y control de fuga}
\label{subsec:split}
PlantVillage incluye varias fotografías de una misma hoja, por lo que una partición aleatoria
puede ubicar imágenes casi idénticas en entrenamiento y en test e inflar las métricas. Para
detectarlas calculamos un hash criptográfico (MD5) de cada archivo y un hash perceptual
(dHash de 64 bits). Encontramos 28 imágenes en 14 grupos de duplicados, todos con etiquetas
consistentes. Luego realizamos una partición estratificada por clase y \emph{agrupada} por
dHash (\texttt{StratifiedGroupKFold}, semilla fija), de modo que las imágenes de un mismo
grupo caen siempre en el mismo subconjunto. El resultado es una partición
71,4/14,3/14,3\,\% (14\,740 imágenes de entrenamiento, 2\,949 de validación y 2\,949 de
test), sin grupos compartidos entre subconjuntos. Esta partición se versiona en el
repositorio y es la misma para todos los experimentos. Como el dHash no es invariante a
rotaciones, este control no elimina la fuga entre fotografías rotadas de una misma hoja; lo
consideramos una limitación del estudio.

\subsection{Preprocesamiento y data augmentation}
\label{subsec:augmentation}
Las imágenes se normalizan con la media y el desvío de ImageNet, que son los usados para
preentrenar los backbones. Validación y test usan siempre un preprocesamiento determinístico
(solo redimensionado). El data augmentation se aplica únicamente en entrenamiento y combina
recortes aleatorios (escala 0,6 - 1,0), volteos horizontales y verticales, rotaciones de hasta
$\pm30^\circ$ ($p=0{,}5$), variaciones moderadas de color (brillo y contraste $\pm30$\,\%,
saturación $\pm20$\,\%, tono $\pm0{,}03$) y desenfoque gaussiano ($p=0{,}2$)
(Fig.~\ref{fig:datos}b). Las hojas no
tienen una orientación canónica, por lo que las transformaciones geométricas son válidas. En
cambio, las variaciones de color se limitaron porque el color es en sí mismo una señal
diagnóstica (clorosis, necrosis).

\subsection{Arquitecturas y estrategia de transferencia}
\label{subsec:arquitecturas}
Elegimos backbones preentrenados en ImageNet con distintos compromisos entre costo y
capacidad: MobileNetV3-Small y MobileNetV3-Large~\cite{howard2019searching},
EfficientNet-B0~\cite{tan2019efficientnet} y, como referencia sin optimización para
dispositivos móviles, ResNet18. Como el entrenamiento se realizó íntegramente en CPU, seguimos
una estrategia en dos etapas:
\begin{enumerate}
  \item \textbf{Linear probe}: con el backbone congelado, extraemos una única vez los
        embeddings de todas las imágenes a $224\times224$ y entrenamos sobre ellos una capa
        lineal. Esta etapa, de pocos segundos por modelo, mide la calidad de las
        representaciones de ImageNet para este dominio y permite comparar las cuatro
        arquitecturas a bajo costo. También evaluamos una entropía cruzada ponderada por la
        frecuencia inversa de cada clase.
  \item \textbf{Fine-tuning parcial}: para las dos arquitecturas con mayor macro-F1 en
        validación en la etapa anterior (EfficientNet-B0 y MobileNetV3-Small, esta última
        además la de menor costo), reemplazamos el clasificador y
        reentrenamos las últimas etapas del backbone (6 bloques en MobileNetV3-Small y 3 en
        EfficientNet-B0, con 1,39\,M y 3,17\,M parámetros entrenables, respectivamente). Las
        BatchNorm de las etapas congeladas se mantienen en modo evaluación para no alterar sus
        estadísticas. Cada arquitectura se entrena con y sin data augmentation, en un diseño
        $2\times2$ en el que todo lo demás (semilla, partición, resolución, optimizador y
        número de épocas) permanece fijo.
\end{enumerate}

\subsection{Protocolo de entrenamiento y métricas}
\label{subsec:entrenamiento}
La Tabla~\ref{tab:config} resume los hiperparámetros de ambas etapas. Para reducir el costo
en CPU, el fine-tuning se realiza a $160\times160$ píxeles. En las dos etapas, la mejor época
se elige por macro-F1 en validación, y el conjunto de test se evalúa una única vez con el
modelo seleccionado. Reportamos accuracy, balanced accuracy y macro-F1, además del F1 por
clase y la matriz de confusión para el análisis de errores. Como métricas de costo
registramos el tiempo de entrenamiento, el número de parámetros y la latencia de inferencia
en CPU.

\begin{table}[t]
\centering
\caption{Configuración de entrenamiento}
\label{tab:config}
\setlength{\tabcolsep}{4pt}
\begin{tabular}{@{}lcc@{}}
\toprule
 & \textbf{Linear probe} & \textbf{Fine-tuning} \\
\midrule
Resolución de entrada      & $224\times224$          & $160\times160$ \\
Capas entrenadas           & capa lineal             & últimas etapas + cabeza \\
Optimizador                & AdamW                   & AdamW \\
Learning rate\textsuperscript{a} & $3\times10^{-3}$  & $10^{-3}$ / $3\times10^{-4}$ \\
Weight decay               & $10^{-3}$               & $10^{-4}$ \\
Scheduler                  & coseno                  & coseno \\
Batch / épocas             & 256 / 60                & 32 / 8 \\
Pérdida                    & CE y CE ponderada       & CE \\
Selección del modelo       & \multicolumn{2}{c}{máximo macro-F1 en validación} \\
\bottomrule
\multicolumn{3}{@{}l}{\footnotesize \textsuperscript{a}En fine-tuning: cabeza / backbone.}
\end{tabular}
\end{table}

\subsection{Explicabilidad con Grad-CAM}
\label{subsec:gradcam}
Para auditar en qué regiones de la imagen se basa el mejor modelo aplicamos
Grad-CAM~\cite{selvaraju2017gradcam} sobre la última capa convolucional del backbone (un mapa
de $5\times5$ a la resolución de entrada), calculado para la clase predicha. Además del
análisis visual, cuantificamos la atención del modelo sobre todo el conjunto de test. Para
cada imagen segmentamos la hoja con un umbral de Otsu sobre el canal de saturación (el fondo
de laboratorio es poco saturado), seguido de limpieza morfológica. Con la máscara $M$
resultante y el mapa Grad-CAM normalizado $L$ calculamos la fracción de la activación que recae 
sobre la hoja:
\begin{equation}
  E_{\text{hoja}} = \frac{\sum_{i,j} L_{ij}\,M_{ij}}{\sum_{i,j} L_{ij}},
\end{equation}
Como un mapa uniforme obtendría un valor igual al área relativa de la hoja, también reportamos 
el cociente entre $E_{\text{hoja}}$ y esa área (mayor que 1 si la atención se concentra en la 
hoja más de lo esperable por azar). Por último, aplicamos el \emph{pointing game}, que mide 
el porcentaje de imágenes cuyo máximo de activación cae dentro de la hoja. Estas métricas 
permiten detectar si el modelo explota el fondo, un sesgo conocido de PlantVillage~\cite{noyan2022uncovering}, 
en lugar de los síntomas.

\subsection{Evaluación de robustez y con datos externos}
\label{subsec:robustez_metodo}
\pendiente{Perturbaciones aplicadas al test (tipos y niveles), dataset externo utilizado,
mapeo de clases y protocolo de evaluación.}
